%------------------------------------------------------------------------------%
\begin{question}
  Encontre uma fórmula explícita para o $n$-ésimo termo da sequência
  \[
    \left(0,\, 3,\, 8,\, 15,\, 24,\, \dots \right)
  \]
\end{question}

%------------------------------------------------------------------------------%
\begin{resolution}{Solução}
  Sequência dos quadrados dos inteiros positivos menos um
  \[
    a_n = n^2-1
  \]
\end{resolution}

%------------------------------------------------------------------------------%
\begin{resolution}{Solução}
  Verificando
  \begin{align*}
    \onslide<+->{a_1 & = (n^2-1)\evalat{n=1}{} = 1^2-1 = 1^2-1 = 0  \\}
    \onslide<+->{a_2 & = (n^2-1)\evalat{n=2}{} = 2^2-1 = 2^2-1 = 3  \\}
    \onslide<+->{a_3 & = (n^2-1)\evalat{n=3}{} = 3^2-1 = 3^2-1 = 8  \\}
    \onslide<+->{a_4 & = (n^2-1)\evalat{n=4}{} = 4^2-1 = 4^2-1 = 15 \\}
    \onslide<+->{a_5 & = (n^2-1)\evalat{n=5}{} = 5^2-1 = 5^2-1 = 24}
  \end{align*}
\end{resolution}

%------------------------------------------------------------------------------%
\begin{resolution}{Solução}
  \emph{Expressão alternativa}
  \[
    a_n = -\frac{1}{3}n^4 +\frac{10}{3}n^3 -\frac{32}{3}n^2 +\frac{50}{3}n -9
  \]
\end{resolution}

%------------------------------------------------------------------------------%
